% 替换位置：摘要中“针对问题二”至“针对问题四”的三段；保留其余摘要
% 这是修改片段，不是完整文档；不要整份直接追加到原文末尾。
针对问题二，建立货箱组批、多点访问、异构机队与共享电池联合调度模型，
以医疗和首批保障时限为硬约束，采用多起点构造与模拟退火邻域搜索求解。
方案使用8架运输无人机和14组电池，通过24架次完成80箱按期交付，
加权延误为0，最晚返航时间为$103.478\,\mathrm{min}$，
总能耗为$69.241\,\mathrm{kWh}$。
本文候选池中，两种目标优先级选中相同方案。

针对问题三，基于三维轨迹、双向链路预算和地形遮挡判定，
将中继服务时间窗口纳入运输排序与机型选择。
方案由24个运输架次和3个中继架次完成，使用2架中继机和3组能源组件，
联合返航时间为$105.037\,\mathrm{min}$，总能耗为$72.435\,\mathrm{kWh}$。
独立回放覆盖爬升、巡航、下降和交接阶段，
经时空分辨率加密，在所述分辨率下未发现通信中断，
但尚未形成连续通信的数学证明。

针对问题四，冻结第三问基准调度，
在组内资源可重新编号、跨组中继完整复制的工作假设下，
将15个服务区合并为8个不可拆分块，利用区间峰值与区间着色核算最小资源，
完整枚举127个两组分区和966个三组分区。
两组方案缺少1架C型运输机和1架中继机；
三组方案缺少1架B型运输机、1架C型运输机、1组B型电池和2架中继机。
两组在资源缺口、配置规模及复制能耗方面更有利。
独立分区的资源需求由组内占用峰值及跨组中继任务复制共同决定。
